% !TEX root = main.tex
\section{시스템 및 신뢰 모델}

\subsection{참여 주체와 권한}
참여자(participant)는 자신이 소유한 객체의 개방값과 소유권
비밀을 관리하고, 증명 클라이언트를 통해 이벤트(event)를 제
출한다. 등록 승인자(entry issuer)는 최초 등록을 승인하며 초
기 값의 적합성과 동일 물량의 중복 등록 방지를 책임진다. 정
책 관리자(policy authority)는 정책 관계와 검증키를 검토하여
등록하고, 정책의 비활성화 및 가공 사용 권한을 관리한다. 검
증 계약(router)과 원장은 증명과 현재 제어 상태를 검증하고,
승인한 전이의 상태 갱신을 원자적으로 실행한다. DPP 검증자
는 외부 문서와 공개 원장으로 제품 주장의 결합을 확인한다.

감사자와 상태 관리자(status authority)는 같은 신뢰 주체
로 둔다. 이 주체는 감사 시작점을 선택하고, 복호화 결과를 해
석하며, 차단을 승인한 범위에 대해 회계 증표와 이전 증표의 상
태를 변경한다. 감사 클라이언트는 감사자가 사용하는 소프트
웨어이고 독립된 신뢰 주체가 아니다. 서로 다른 운영 역할을
한 주체가 맡더라도 각 역할의 권한 검사는 계속 적용한다.

소유권은 공개 트랜잭션(transaction) 주소와 구분한다. 소
유권 비밀 \(sk\)에 대해 \(o=\) OwnerKey\((sk)\)를 정의하고, 객체를
소비하는 증명은 해당 비밀의 지식을 요구한다. 공개 제출자가
소유자와 동일한지, 중계자를 허용하는지와 같은 인증 방식은
배포 설정에서 정한다. 초기 등록(Entry)의 생성자는 최초 회계
증표의 소유자이고, 병합(Merge)과 분할(Split)은 소유권을 유
지한다. 가공(Process)의 모든 입력과 출력은 공정 수행 주체의
소유다. 조직 간 이전에서는 이전 이벤트(Transfer)로 이전 증
표를 생성하고 수신자가 인수(Proceed)로 이를 해결하는 경로
를 사용한다. 이에 따라 출력 회계 증표를 생성하는 주체는 그
회계 증표의 미래 소비 식별자(nullifier)를 계산할 소유권 비밀
을 알고, 송신자는 자신이 생성한 이전 증표의 개방값으로 그 증

% Source PDF page 5, column 1
표의 미래 소비 식별자(resolution nullifier)를 계산한다.

\subsection{감사 복호화 권한}
감사 기록의 암호화와 그 복호화 권한의 확보에 서로 다른
키를 사용한다. 일반 공개키 암호의 감사 데이터 키쌍을
\((PK_{B},SK_{B})\)로 두고, 본문의 \(pk_{audit}=PK_{B}\), \(sk_{audit}=\)
\(SK_{B}\)로 표기한다. 참여자는 \(PK_{B}\)로 감사 평문을 암호화한
다. 위원회(committee)는 별도의 임계 암호화(threshold en-
cryption) 공개키 \(PK_{A}\)와 구성원별 비밀 조각을 관리한다. 신
뢰된 초기 설정은 데이터 비밀키를 다음과 같이 보호한다.

\begin{equation}
    W_B=\operatorname{TE.Enc}(PK_A,\operatorname{EncodeSecret}(SK_B),\eta_A;\mathrm{ctx}_{\mathrm{wrap}}).\tag{1}
\end{equation}
여기서 \(W_{B}\)는 데이터 비밀키를 보호하는 암호문이고, \(η_{A}\)는 그
암호화 난수이다. 외부 승인 절차는 \(N\)명 중 \(K\)명 이상의 유효한
협력으로 \(W_{B}\)의 평문을 복구하여 승인된 감사자에게 \(SK_{B}\)를
인증된 비밀 경로로 전달한다. \(1≤K≤N\)이며, 위원회의 원
래 비밀 조각을 감사자에게 전달하는 것을 요구하지 않는다.

승인된 설정은 배포와 키 식별자, \(PK_{A}\), \(PK_{B}\), \(W_{B}\)의 대응
을 인증해야 한다. 초기 설정은 \(W_{B}\)의 평문이 승인된 \(PK_{B}\)에
대응하는 유효한 비밀키인지 확인할 책임이 있다. 생성 주체와
비밀 사본의 보관·폐기 책임도 설정의 신뢰 경계에 포함된다.
초기 보유자가 \(SK_{B}\)의 사본을 유지한다면 그 주체도 기록을 복
호화할 수 있으므로, 임계 암호화만으로 모든 키 접근이 위원회
협력을 통해서만 가능하다고 보지 않는다. \(ctx_{wrap}\)은 데이터 키
보호라는 목적과 배포·키 식별을 \(PK_{B}\)에 결합한다. 기록 암호
화의 문맥은 별도로 전이와 슬롯을 결합하며, 두 결합의 정규 인
코딩과 인증 방법은 배포 설정에서 고정한다.

감사 내부의 복호화 인터페이스는

\begin{equation}
    m\leftarrow\operatorname{AuditDec}(sk_{\mathrm{audit}},ct;\mathrm{ctx})\tag{2}
\end{equation}
이며 실패 시 오류를 반환한다. 감사의 비밀 로컬 입력은 승인
된 키 식별자와 유효한 \(SK_{B}\)이며, 키 식별·정규형·키쌍 확인에
서 발견된 불일치는 거부한다. 데이터 비밀키는 원장 트랜잭션
이나 공개 감사 결과에 포함하지 않는다. 이 권한 확보 이후 감
사자는 해당 키 영역의 기록을 로컬에서 복호화한다. 방문한 암
호문마다 위원회에 원격 부분 복호화를 요청하는 절차는 감사
내부에 포함하지 않는다. 위원회 협력 중 부분 복호화를 사용한
다면 그 대상은 \(W_{B}\)이다. 승인·복구·전달과 잘못된 응답의 처
리는 외부 절차의 책임이며, 이 논문은 유효한 데이터 키가 확보
되었다는 경계에서 추적과 차단을 분석한다. 식 (1)는 두 키의
역할을 정하며 특정 암호 방식의 보안이나 키 전달 절차의 검증
완료를 뜻하지 않는다.

해당 키를 확보한 감사자는 그 키로 보호된 다른 암호문도 복
호화할 수 있다. 감사 시작점과 시점은 실행할 계산의 범위를 정
하지만 키 사용을 그 암호문 집합으로 제한하는 암호학적 장치
는 아니다. 감사 자료와 수령 키의 취급은 신뢰된 감사자의 책
임이다. 기본 실행은 배포에 고정된 하나의 데이터 키를 사용한
다. 키 교체를 지원하는 확장에는 과거 기록의 키 식별과 접근·
실패 처리가 추가로 필요하며, 키 교체나 승인 철회가 이미 전달
한 \(SK_{B}\)의 복호화 능력을 회수하지는 않는다.

\subsection{공격자와 외부 가정}
일반 참여자, 증명자, 제출자와 DPP 게시자는 공모하거나 잘못
된 증인·증명·문서를 제출할 수 있다. 비인가 공격자는 감사 데

% Source PDF page 5, column 2
이터 비밀키 \(SK_{B}\)를 얻지 못하고, 위원회 비밀 조각은 \(K\)개 미
만만 얻는다고 가정한다. 초기 설정 사본과 키 전달 채널을 통
한 비인가 노출도 배제하며, 그 경로로 유출된 \(SK_{B}\)를 임계 조
건이 보호한다고 주장하지 않는다. 승인된 감사자의 키 접근은
비인가 노출과 구분하며, 그 감사자에 대한 감사 정보의 비밀성
은 요구하지 않는다. 등록 승인자의 초기 판단, 정책 관리자의
정책·권한 등록, 상태 관리자의 시작점 선택과 올바른 식별자 집
행을 신뢰한다.

원장은 확정된 실행 순서, 변경되지 않는 수락 기록, 인증된
상태 읽기와 원자적 트랜잭션을 제공한다고 가정한다. 외부 색
인 제공자나 원격 프로시저 호출(remote procedure call, RPC)
의 응답을 인증 없이 신뢰하지 않는다. 이전 증표 개방값은 송
신자가 생성하여 보관하고, 수신자에게 인증된 비밀 채널로 동
일한 값을 전달한다고 가정한다. 이 전달의 실패 복구와 가용성
을 원장 관계만으로 보장하지 않는다.

질량, 재활용 크레딧, 초기 탄소와 운송·공정 증분은 외부에
서 제공된다. 증명은 이 값의 범위와 전이 관계를 검사하며 실
제 측정의 정확성, 물리적 공정의 수행 또는 문서와 실제 제품의
동일성을 판정하지 않는다.
